\section{Arithmetic realizations of the identity line}
\label{sec:weighted-cartier-criterion}\label{sec:prime-tests}

The coefficient line has a second realization in characteristic $p$.
The resulting criterion is unconditional, but its hypothesis concerns
reductions at primes.  We shall keep this hypothesis separate from the
complex period relation.

\subsection{From a finite polynomial to a de Rham line}
For $s\in\{1/6,1/4,1/3,1/2\}$ and $p>3$, put
\[
 U_p(x)=\sum_{n<p}\frac{(s)_n(1-s)_n}{(n!)^2}x^n,
 \qquad T_p(z)=\sum_{n<p}c_s(n)z^n,
 \qquad W_p(z)=(A+B\thetaop)T_p(z).
\]
Thus $W_p$ is the canonical truncation of the proposed identity.
Let $z=4x(1-x)$, write $\Omega=dX/y$ on the standard elliptic
model, and put
\[
 Q(x)=\frac{2x(1-x)}{1-2x},\qquad
 \nu_{A,B}=A\Omega+BQ(x)\nabla_x\Omega\in H_x,
 \qquad H_x=H^1_{\mathrm{dR}}(E_{s,x}).
\]
We exclude finitely many primes at which the parameter, this basis,
or a specified nonzero coefficient fails to be integral or a unit.
All subsequent density statements are insensitive to these exclusions.

The precise unconditional polynomial input from
\cite[\S7, equation (7.4) and Lemma 7.2]{BRSTT} is
\begin{equation}\label{eq:finite-clausen-p2}
 T_p(4x(1-x))\equiv U_p(x)^2\pmod{p^2\Z_p[x]}.
\end{equation}
Differentiating this identity gives, after specialization away from
$1-2x=0$,
\begin{equation}\label{eq:cartier-weighted-factorization}
 W_p(z)\equiv U_p(x)\bigl(AU_p(x)+BQ(x)U_p'(x)\bigr)
                                                   \pmod{p^2}.
\end{equation}
In particular, this assertion makes no CM assumption.  The source's
CM supercongruence is a different result; its hypothesis and first-order
term are recorded in Appendix~\ref{subsec:precision-versus-transfer}.

\begin{lemma}[The Cartier realization]
\label{lem:cartier-hypergeometric-bridge}
Let $k$ be a perfect residue field at a good place, with
$x(1-x)(1-2x)\ne0$.  The reduction is ordinary if and only if
$U_p(x)\ne0$.  In that case, for $B\ne0$ in $k$,
\begin{equation}\label{eq:rev:cartier-line}
 W_p(z)=0
 \quad\Longleftrightarrow\quad
 k\nu_{A,B}=\ker(\operatorname{Cartier})
             =\operatorname{im}(F\bmod p).
\end{equation}
The displayed line is the unique Frobenius-stable complement to the
Hodge line $F^1H_x$.  These statements hold over every perfect residue
field, including quadratic extensions of $\mathbf F_p$.
\end{lemma}
\begin{proof}
The Hasse invariant in our differential normalization is
$(-6)^{(p-1)/2}U_p(x)$.  More generally, Cartier evaluation of
$\nu_{A,B}$ is, up to this fixed nonzero factor and the Frobenius
twist, the linear functional
\[
             (A,B)\longmapsto AU_p(x)+BQ(x)U_p'(x).
\]
The normalization and its compatibility with Gauss--Manin differentiation
are proved in Appendix~\ref{rev:cartier-calculation}.  Conceptually,
this calculation identifies the top row of the diagram
\[
\begin{tikzcd}[column sep=large]
 k^2 \arrow[r,"{(A,B)\mapsto\nu_{A,B}}"]
       \arrow[dr,swap,"(AU_p+BQU_p')^{1/p}"]
   & H_x \arrow[d,"\operatorname{Cartier}"]\\
   & H^0(E_{s,x},\Omega^1)^{(p)},
\end{tikzcd}
\]
after rescaling its target basis by the fixed nonzero normalization.
The diagonal map is semilinear, as is Cartier.
The conjugate-filtration exact sequence is
\begin{equation}\label{eq:rev:conjugate-sequence}
 0\longrightarrow H^1(E_{s,x},\mathcal O)^{(p)}
 \xrightarrow{F}H_x
 \xrightarrow{\operatorname{Cartier}}
 H^0(E_{s,x},\Omega^1)^{(p)}\longrightarrow0.
\end{equation}
Its first image is a line.  Since the Gauss--Manin coefficient of
$\Xi$ in $Q\nabla_x\Omega$ is nonzero, $\nu_{A,B}$ spans a
complement to $F^1H_x$ when $B\ne0$.  At an ordinary fiber $U_p(x)$
is a unit, so \eqref{eq:cartier-weighted-factorization} identifies
its vanishing with the kernel in \eqref{eq:rev:conjugate-sequence}.
Finally $F\bmod p$ kills the Hodge line and has rank one.  A stable
complement must therefore equal its image; at an ordinary fiber that
image is indeed a complement.
\end{proof}

At a supersingular fiber the first factor in
\eqref{eq:cartier-weighted-factorization} vanishes for every pair
$(A,B)$.  Such zeros do not assert stability of a complementary line.
This distinction determines the density argument.

\subsection{Algebraization and the density of stable complements}
\label{subsec:B-tang-packet}\label{subsec:B-hodge-complements}
For a number field $K$ and a set $M$ of finite places, define
\begin{equation}\label{eq:rev:weighted-density}
 d_K(M;X)=
 \frac{\displaystyle\sum_{\substack{v\in M\\p_v\le X}}
 [K_v:\Q_{p_v}]\frac{\log p_v}{p_v-1}}
 {[K:\Q]\displaystyle\sum_{p\le X}\frac{\log p}{p-1}},
 \qquad
 \underline d_K(M)=\liminf_Xd_K(M;X),\quad
 \overline d_K(M)=\limsup_Xd_K(M;X).
\end{equation}
Here $p_v$ is the rational residue characteristic, not the norm of
$v$.  Finite sets have density zero, and
\begin{equation}\label{eq:B-density-complement}
               \underline d_K(M)=1-\overline d_K(M^c).
\end{equation}

\begin{theorem}[Tang's algebraization inequality]
\label{thm:B-Tang-external}
Let $X/K$ be a geometrically irreducible quasi-projective variety of
dimension $N$, let $P\in X(K)$, and let $\mathcal D\subset T_X$
be an involutive subbundle of rank $d$.  Suppose that its smooth
formal leaf through $P$ is Zariski dense.  For every embedding
$\sigma:K\hookrightarrow\mathbb C$, assume a holomorphic map
$\gamma_\sigma:\mathbb C^d\to X_\sigma(\mathbb C)$ takes zero
to $P_\sigma$ and induces that leaf germ biholomorphically.
Suppose these maps have finite Nevanlinna orders $\rho_\sigma$,
and put $\rho=\max_\sigma\rho_\sigma$.

Spread these data over $\mathcal O_K[1/n]$.  Let $M_{\mathrm{bad}}$
be the places where the reduced distribution is not closed under
restricted $p_v$th powers, and put
$\alpha_{\mathrm{bad}}=\overline d_K(M_{\mathrm{bad}})$.
Then either $N=d$, or
\begin{equation}\label{eq:B-Tang-inequality}
                  1\le\frac{N}{N-d}\rho\alpha_{\mathrm{bad}}.
\end{equation}
\end{theorem}
This is the quantitative form of
\cite[\S\S5.1.1--5.1.4 and Theorem 5.1.5]{Tang}.  The order is
computed using an ample Hermitian line bundle on a projective
compactification; the needed normalization is given in
Appendix~\ref{subsec:B-growth-sharp}.  The use of an upper density
for the bad places is essential.

\begin{lemma}[Naturality of the density]
\label{lem:B-density-transport}
For a finite extension $K'/K$, let $M'$ consist of all places above
$M$.  Then $d_{K'}(M';X)=d_K(M;X)$.  If $M$ contains every place
above a set $S$ of rational primes of lower natural density $\delta$,
apart from finitely many exceptions, then
$\underline d_K(M)\ge\delta$.
\end{lemma}
\begin{proof}
The first assertion is the local degree formula
\[
 \sum_{w\mid v}[K'_w:\Q_{p_v}]
            =[K':K][K_v:\Q_{p_v}].
\]
For the second, sum the local degrees over each rational prime and
apply partial summation to the positive decreasing weight
$w(t)=\log t/(t-1)$.  The bound
$\#\{p\le t:p\in S\}\ge(\delta-\epsilon)\pi(t)$, valid for
large $t$, gives the corresponding weighted bound up to a constant.
Division by $\sum_{p\le X}w(p)\sim\log X$ and then
$\epsilon\downarrow0$ proves the assertion.
\end{proof}
Thus all places over an admitted rational prime must be retained.
For instance, an inert prime in a quadratic field contributes its
full weight here, although its prime ideal has norm $p^2$ and the
set of such prime ideals has Dirichlet density zero.

\begin{theorem}[An algebraic complement cannot be too often stable]
\label{thm:B-complement-general}\label{lem:complement-density-bound}
\label{lem:B-full-complement-dense}
Let $A/K$ be an abelian variety of positive dimension and let
$W\subset H^1_{\mathrm{dR}}(A/K)$ be a $K$-linear complement to
$F^1H^1_{\mathrm{dR}}(A/K)$.  Let $M_W$ be the good places where
its reduction is stable under crystalline Frobenius modulo the
residue characteristic.  Then
\begin{equation}\label{eq:B-complement-density}
                         \underline d_K(M_W)\le\frac34.
\end{equation}
No simplicity or non-CM assumption is needed.
\end{theorem}
\begin{proof}
Let $g=\dim A$ and let $G=\mathcal E(A^\vee)$ be the universal
vector extension.  Its defining sequence realizes the Hodge
filtration as tangent spaces:
\begin{equation}\label{eq:rev:universal-extension-diagram}
\begin{tikzcd}[column sep=small]
 0\arrow[r]&V\arrow[r]\arrow[d,phantom,"\scriptstyle\operatorname{Lie}" description]
 &G\arrow[r]\arrow[d,phantom,"\scriptstyle\operatorname{Lie}" description]
 &A^\vee\arrow[r]\arrow[d,phantom,"\scriptstyle\operatorname{Lie}" description]&0\\
 0\arrow[r]&F^1H^1_{\mathrm{dR}}(A)\arrow[r]
 &H^1_{\mathrm{dR}}(A)\arrow[r]&H^1(A,\mathcal O_A)\arrow[r]&0.
\end{tikzcd}
\end{equation}
At good places this identification is integral, and the restricted
$p$th power on $\operatorname{Lie}G$ is crystalline Frobenius
modulo $p$; these are the comparison facts used in
\cite[proof of Theorem 2.3.5 and Lemma 5.2.6]{Tang}.

The exponential leaf of $W$ is Zariski dense in $G$.  Indeed, let
$H$ be its connected algebraic closure and $U=H\cap V$.
Surjectivity on tangent spaces shows that $H\to A^\vee$ is
surjective.  A subgroup of a vector group in characteristic zero
is a vector group, so $H/U\simeq A^\vee$ splits the quotient
extension $G/U$.  On extension classes this would make zero the
restriction
\[
 (V/U)^*\lhook\joinrel\longrightarrow V^*
       \xrightarrow{\ \sim\ }H^1(A^\vee,\mathcal O_{A^\vee}).
\]
The last map is the universal extension class.  Its injectivity
forces $U=V$, and hence $H=G$.

Translation of $W$ now defines a rank-$g$ involutive distribution
on the $2g$-dimensional group $G$.  At every complex embedding
its exponential uniformization has order at most two, by the
commutative-group estimate in
\cite[proof of Proposition 2.3.4, footnote 20]{Tang}.
Restricted-power stability of this distribution is precisely the
condition defining $M_W$.  Theorem~\ref{thm:B-Tang-external}
therefore gives
\[
 1\le\frac{2g}{g}\rho\,\overline d_K(M_W^c)
                \le4\overline d_K(M_W^c).
\]
Equation~\eqref{eq:B-density-complement} finishes the proof.
\end{proof}

\subsection{The arithmetic classification}
\label{subsec:B-Tang-application}
We need one Galois-theoretic input that counts rational primes,
including those of higher residue degree.

\begin{theorem}[Supersingular density for $\Q$-curves]
\label{thm:B-qcurve-density}\label{lem:qcurve-density}
Let $E$ over a number field be geometrically isogenous to every
Galois conjugate.  Outside finitely many rational primes,
supersingularity is independent of the place above the rational
prime.  Its rational-prime set has natural density zero if $E$ is
non-CM, and density $1/2$ if $E$ is CM.
\end{theorem}
\begin{proof}
For non-CM $E$, the one-dimensional rational spaces of isogenies
between its conjugates assemble its Tate representation into a
projective representation of $G_\Q$.  Serre's open image theorem
\cite{Serre} makes every sufficiently large residual image contain
$\operatorname{PSL}_2(\mathbf F_\lambda)$.  At a supersingular
rational prime its projective Frobenius has bounded order, with a
bound depending only on a fixed Galois field of definition.
Elements of bounded order have proportion $O(1/\lambda)$ in these
images.  Chebotarev, followed by $\lambda\to\infty$, gives
density zero.  Appendix~\ref{rev:qcurve-proof} proves both the
projective construction and the uniform order bound, so this
argument includes primes of higher residue degree.

In the CM case, Deuring's reduction theorem \cite{Deuring} says
that good rational primes unramified in the CM field are
supersingular exactly when inert in that imaginary quadratic
field.  Their density is $1/2$ by Chebotarev.  Specialization of
the isogenies shows independence of the place in both cases.
\end{proof}

\begin{theorem}[Weighted-density CM criterion]
\label{thm:weighted-density-cm}
Let $z\in\Q\setminus\{0,1\}$, let
$x=(1-\sqrt{1-z})/2$, and let $(A,B)\in\Q^2\setminus\{0\}$.
If
\[
       \underline\delta\{p:W_p(z)\equiv0\pmod p\}>\frac34,
\]
then $E_{s,x}$ has CM.  In particular it suffices that the
congruence hold at all but finitely many ordinary rational primes.
\end{theorem}
\begin{proof}
Suppose $E_{s,x}$ is non-CM.  The standard modular correspondences
make it a $\Q$-curve over $K=\Q(x)$, so its supersingular rational
primes have density zero by Theorem~\ref{thm:B-qcurve-density}.
If $B=0$, \eqref{eq:cartier-weighted-factorization} excludes every
ordinary zero outside a finite set.  If $B\ne0$, the fixed line
$K\nu_{A,B}$ complements the Hodge line.  At every admitted
ordinary rational prime, Lemma~\ref{lem:cartier-hypergeometric-bridge}
makes it Frobenius-stable at \emph{every} place of $K$ above that
prime.  Indeed $W_p(z)$ is rational, and the same polynomial
calculation applies at both conjugate parameters and over their
residue fields.  Lemma~\ref{lem:B-density-transport} consequently
gives lower weighted-place density greater than $3/4$, contrary to
Theorem~\ref{thm:B-complement-general}.  If all but finitely many
ordinary primes are admitted, the non-CM supposition gives density
one and the same contradiction.
\end{proof}

\begin{proposition}[The CM density dichotomy]
\label{prop:cm-weighted-slope}
Suppose $E_{s,x}$ has CM, and denote its complementary CM eigenline
by $L_x^{\mathrm{CM}}$, as in \eqref{eq:cm-canonical-splitting}.
For every fixed nonzero rational pair $(A,B)$,
\[
 \delta\{p:W_p(z)=0\bmod p\}=
 \begin{cases}
 1,&K\nu_{A,B}=L_x^{\mathrm{CM}},\\
 1/2,&K\nu_{A,B}\ne L_x^{\mathrm{CM}}.
 \end{cases}
\]
In the first case the zero set is cofinite; in the second its
ordinary part is finite.
\end{proposition}
\begin{proof}
The companion eigenline descends to $K$ because Galois preserves
the pair of CM eigenspaces and fixes the Hodge line.  At an
ordinary prime $p$, the CM field splits.  Its two eigenvalues lie
in $\Q_p$ and are fixed by coefficient Frobenius; functoriality
therefore makes crystalline Frobenius preserve both eigenlines.
Modulo $p$ it kills the Hodge line and has rank one, whence
\[
                  \operatorname{im}(F\bmod p)
                          =L_x^{\mathrm{CM}}\bmod p.
\]
The equality descends from a field defining the CM action.  Thus
the Cartier criterion identifies ordinary weighted zeros with
equality of the reductions of $K\nu_{A,B}$ and
$L_x^{\mathrm{CM}}$.  Distinct algebraic lines have distinct
reductions outside finitely many places: their exterior product
is nonzero and has only finitely many prime divisors.  At every
supersingular prime the Hasse factor $U_p(x)$ vanishes, so every
pair gives a weighted zero.  Deuring's density $1/2$ proves all
assertions, including the case $B=0$.
\end{proof}

\begin{theorem}[Arithmetic classification of primitive rows]
\label{thm:weighted-row-classification}
Fix $\ell\in\{1,2,3,4\}$.  Among integer triples with
$\gcd(A,B)=1$, $|C|\ge R_\ell$ and
$C\ne R_\ell$, the following conditions are equivalent:
\begin{enumerate}
\item the lower natural density of
 $\{p:W_p(R_\ell/C)\equiv0\pmod p\}$ is greater than $3/4$;
\item this set of primes is cofinite;
\item up to simultaneous sign of $A,B$, the triple is one of the
 standard rows at level $\ell$ in Section~\ref{sec:standard}.
\end{enumerate}
\end{theorem}
\begin{proof}
The first condition forces CM.  The finite classification in
Section~\ref{sec:enumeration} leaves the tabulated parameters.
At each parameter Proposition~\ref{prop:cm-weighted-slope}
selects exactly the complementary CM eigenline, which is the
complex identity line by Proposition~\ref{prop:cm-matrix-certificate}
and projective uniqueness.  Thus the first condition gives the
third, and the third gives the second.  The remaining implication
is immediate.
\end{proof}
Consequently, on the original convergence domain, CM forcing is
equivalent to the assertion that every complex identity has
weighted-zero lower density greater than $3/4$.  The arithmetic
criterion describes the missing comparison statement; it does
not prove that statement.

\subsection{An exact comparison at one ordinary prime}
\label{sec:one-prime-c}
There is a second arithmetic endpoint: exact preservation of the
characteristic-zero Hodge line by a single ordinary Frobenius
already characterizes CM.

\begin{theorem}[One-prime CM criterion]\label{thm:one-prime-c}
Let $E/K$ be an elliptic curve over a number field and let $v\mid p$
be a place of good ordinary reduction.  Write $L=K_v$, let
$k=\mathbf F_{p^f}$ be its residue field, and put
\[
 K_0=\operatorname{Frac}W(k),\qquad
 D=H^1_{\mathrm{cris}}(E_k/W(k))[1/p],\qquad\Phi=\varphi^f.
\]
Through crystalline--de Rham comparison, extend the $K_0$-linear
endomorphism $\Phi$ to $H=H^1_{\mathrm{dR}}(E/L)$.  Then
\[
 E\text{ has geometric CM}
 \quad\Longleftrightarrow\quad
                    \Phi(F^1H)\subseteq F^1H.
\]
This formulation allows $L/K_0$ to be ramified and selects no
Frobenius lift on $L$.
\end{theorem}
\begin{proof}
We use the $p$-adic Lefschetz theorem in the following precise form
\cite[Theorem 4.24]{MaulikPoonen}.  Let $\mathcal X$ be a smooth
proper $p$-adic formal scheme over the valuation ring of a complete
discretely valued characteristic-zero field with perfect residue
field.  A line bundle $\mathcal L_0$ on its special fiber has some
tensor power $\mathcal L_0^{\otimes p^m}$ lifting to
$\mathcal X$ if and only if its crystalline Chern class, transported
to de Rham cohomology, belongs to $F^1H^2_{\mathrm{dR}}(\mathcal X)$.
Ramification is permitted.

Let $\mathcal E/\mathcal O_L$ be the good model and let
$\pi_q:E_k\to E_k$ be $q$-power Frobenius, $q=p^f$.
Its graph is the divisor
$\Gamma=(\pi_q,-1)^*\{0\}$ on $E_k\times E_k$.
The polarization identifies the mixed K\"unneth component of
$c_1(\mathcal O(\Gamma))$ with the crystalline action of
$\pi_q$, namely $\Phi$.  Under this identification the obstruction
to the graph class lying in $F^1$ is exactly the map
\begin{equation}\label{eq:rev:frobenius-obstruction}
                 F^1H\xrightarrow{\ \Phi\ }H
                           \longrightarrow H/F^1H.
\end{equation}
Indeed the polarization gives $F^1H\simeq(H/F^1H)^*$ after
trivializing the one-dimensional determinant.  Projection of the
correspondence tensor to $(H/F^1H)^{\otimes2}$ is consequently
the same obstruction.  The two unmixed K\"unneth components
are divisor classes on the factors and already lie in $F^1$.

If \eqref{eq:rev:frobenius-obstruction} vanishes, the quoted
Lefschetz theorem lifts $\mathcal O(\Gamma)^{\otimes p^m}$
to the formal completion of $\mathcal E\times\mathcal E$.
Formal GAGA algebraizes the lift to a line bundle $\mathcal L$
\cite[Example 4.5(b)]{MaulikPoonen}.  For a line bundle on a
product of principally polarized elliptic curves, the negative
off-diagonal block of its associated homomorphism to the dual
is an endomorphism.  This construction is additive in the line
bundle and commutes with base change:
\[
\begin{tikzcd}[column sep=large]
 \operatorname{Pic}(\mathcal E\times\mathcal E)
      \arrow[r,"\text{correspondence}"]\arrow[d,"\mathrm{sp}"']
   &\operatorname{End}(\mathcal E)\arrow[d,"\mathrm{sp}"]\\
 \operatorname{Pic}(E_k\times E_k)
      \arrow[r,"\text{correspondence}"']
   &\operatorname{End}(E_k).
\end{tikzcd}
\]
For the graph divisor its value is $\pi_q$; hence $\mathcal L$
gives an endomorphism $u$ specializing to $[p^m]\pi_q$.
Ordinary Frobenius has eigenvalue valuations $0,f$, so this
specialization is nonscalar.  A scalar generic endomorphism would
specialize to multiplication by an integer, which has equal
valuations on both eigenvalues.  Thus $u$ is nonscalar.  Geometric
homomorphism groups of abelian varieties are unchanged by extension
of algebraically closed fields; consequently $E_{\overline K}$
has a nonscalar endomorphism and has CM.

Conversely, suppose $E$ has CM.  After a finite local extension,
all its geometric endomorphisms extend to the good model.
Specialization embeds its rational geometric endomorphism algebra
in that of the ordinary special fiber.  Both are imaginary
quadratic fields, so the embedding is an isomorphism.  Thus
$\pi_q$ is the specialization of a characteristic-zero
quasiendomorphism.  Its action preserves the Hodge filtration,
and crystalline--de Rham comparison is functorial for this action
\cite[Theorem 4.21]{MaulikPoonen}.  The vanishing of
\eqref{eq:rev:frobenius-obstruction} descends from that extension
to $L$.
\end{proof}

\begin{remark}\label{rem:one-prime-scope-c}
The lifting argument applies to any special-fiber endomorphism
whose crystalline action preserves the characteristic-zero Hodge
filtration.  Ordinarity guarantees that Frobenius is nonscalar;
at supersingular reduction a power of Frobenius can be scalar
even for a non-CM characteristic-zero curve.  Furthermore,
stability of the complementary line modulo $p$ in
Lemma~\ref{lem:cartier-hypergeometric-bridge} is a different
condition from exact stability of the Hodge line here.
Appendix~\ref{prop:one-prime-precision-c} shows that no fixed
finite precision substitutes for exactness.  No implication
from the original complex scalar identity to this exact
Frobenius condition is established.
\end{remark}
